\documentclass[11pt]{article}
\usepackage[margin=0.95in]{geometry}
\usepackage{amsmath,amssymb,lmodern}
\usepackage[T1]{fontenc}
\usepackage{needspace}
\usepackage[colorlinks=true,linkcolor=blue,urlcolor=blue]{hyperref}
\hypersetup{pdftitle={Positivity and counterexamples for rising-factorial Toeplitz determinants},pdfauthor={Henry Zweiman}}
\setlength{\emergencystretch}{3em}
\setlength{\parskip}{0.3em}
\setcounter{secnumdepth}{0}
\providecommand{\tightlist}{\setlength{\itemsep}{0pt}\setlength{\parskip}{0pt}}
\title{Positivity and counterexamples for rising-factorial Toeplitz determinants}
\author{Henry Zweiman}
\date{September 22, 2026}
\begin{document}
\maketitle
\begin{center}\small Research preprint, version 1.0.\\
Prepared with OpenAI Codex; not peer reviewed.\end{center}
\Needspace{9\baselineskip}\section{Abstract}\label{abstract}

We examine the conjectures other than Conjecture 3 in
Karp\textquotesingle s 2012 paper on rising-factorial Toeplitz
determinants. A positive expansion in two-row Toeplitz minors proves
Conjecture 1 for all positive shift parameters. A finite-difference and
Cauchy-\/-Binet formula proves coefficient nonnegativity for arbitrary
determinant order under the corresponding Polya frequency hypothesis,
with exact criteria for nonvanishing, degree, and strict positivity. It
also identifies the degeneracies excluded by a corrected form of
Conjecture 4. The stability assertions in Conjectures 2 and 5 are false:
a positive strictly log-concave sequence of length nineteen gives a
negative Hurwitz determinant. Conjecture 6 is false even for the
binomial sequence of degree eleven; an order-three determinant has a
cubic factor with negative discriminant. All counterexamples have exact
algebraic certificates. This is a separate manuscript from the
author\textquotesingle s paper on Conjecture 3.

\Needspace{9\baselineskip}\section{1. Definitions and results}\label{definitions-and-results}

Write \((x)_0=1\) and \((x)_k=x(x+1)\cdots(x+k-1)\) for \(k\ge1\). Let
\(f=(f_0,f_1,\ldots)\) be a nonnegative sequence, and work with formal
power series

\[
\Phi_f(x,z)=\sum_{k\ge0} f_k(x)_k\frac{z^k}{k!}.
\tag{1}
\]

Only \(f_0,\ldots,f_n\) affect any coefficient of \(z^n\) below. A
finite input is extended by zero. Following the explicit multinomial
definition in {[}K12{]}, we use the normalization

\[
\begin{aligned}
Q_n^{\alpha,\beta}(x)
&=n![z^n]\bigl(\Phi_f(x+\alpha,z)\Phi_f(x+\beta,z)\\
&\hspace{37mm}-\Phi_f(x+\alpha+\beta,z)\Phi_f(x,z)\bigr),\\
P_n^r(x)&=n![z^n]\det\bigl[\Phi_f(x+j-i,z)\bigr]_{i,j=0}^{r-1}.
\end{aligned}
\tag{2}
\]

Thus \(P_n^2(x)=Q_n^{1,1}(x-1)\). Multiplication of either definition by
a positive normalization constant has no effect on the assertions about
zeros or coefficient signs.

A sequence is \(PF_r\) if every minor of order at most \(r\) of its
Toeplitz matrix \([f_{j-i}]_{i,j\ge0}\) is nonnegative, where \(f_k=0\)
for \(k<0\). The notation \(PF_\infty\) requires this in every order. In
particular, \(PF_2\) is equivalent to log-concavity with no internal
zeros. Hurwitz stability means that every zero has strictly negative
real part. The zero polynomial is treated separately throughout.

The disposition of the remaining conjectures is as follows.

\begin{itemize}
\tightlist
\item
  \textbf{Conjecture 1:} proved for every \(\alpha,\beta>0\) in Theorem
  1.
\item
  \textbf{Conjecture 2:} disproved by the strictly log-concave positive
  input (15), with \(\alpha=\beta=1\).
\item
  \textbf{Conjecture 4:} the unconditional degree and strict-positivity
  assertion fails for degenerate inputs. Theorem 2 proves coefficient
  nonnegativity and gives necessary and sufficient minor conditions for
  the asserted degree and strict positivity.
\item
  \textbf{Conjecture 5:} disproved already for \(r=2\) by the same
  positive input (15).
\item
  \textbf{Conjecture 6:} disproved for \(r=3\), \(n=11\), and
  \(f_k=\binom{11}{k}\).
\end{itemize}

The numbers in this list refer to {[}K12{]}. These are dispositions of
the stated assertions, not a claim that all of them are true. Section 6
explains the separate scope of Conjecture 3. No historical priority or
exhaustive literature determination is asserted.

\Needspace{9\baselineskip}\section{2. Arbitrary positive shifts: proof of Conjecture
1}\label{arbitrary-positive-shifts-proof-of-conjecture-1}

For \(j\ge0\), define the shifted series

\[
G_j(x,z)=\sum_{a\ge0}f_{a+j}(x)_a\frac{z^a}{a!}.
\]

\Needspace{5\baselineskip}\textbf{Theorem 1.} Suppose \(f_0,\ldots,f_n\) is nonnegative,
log-concave, and has no internal zeros, and \(\alpha,\beta>0\). Then
\(Q_n^{\alpha,\beta}\) has nonnegative coefficients and degree at most
\(n-2\) when nonzero. If \(n\ge3\) and

\[
f_k^2>f_{k-1}f_{k+1}\qquad(1\le k<n),
\tag{3}
\]

then its degree is exactly \(n-2\), and every coefficient in degrees
\(0,\ldots,n-2\) is strictly positive.

\textbf{Proof.} The rising-factorial Vandermonde identity gives

\[
\Phi_f(x+\alpha,z)=\sum_{j\ge0}\frac{(\alpha)_jz^j}{j!}G_j(x,z).
\]

Expanding the first shift in the two products defining \(Q\) therefore
gives

\[
\begin{aligned}
&\Phi_f(x+\alpha,z)\Phi_f(x+\beta,z)
-\Phi_f(x+\alpha+\beta,z)\Phi_f(x,z)\\
&\quad=\sum_{j\ge1}\frac{(\alpha)_jz^j}{j!}
\bigl(G_j(x,z)G_0(x+\beta,z)
-G_j(x+\beta,z)G_0(x,z)\bigr).
\end{aligned}
\]

Pair the summands with indices \((a,b)\) and \((b,a)\) in the last
product difference. Set

\[
\begin{aligned}
M_{a,b}^{(j)}&=f_{a+j}f_b-f_af_{b+j},\\
B_{a,b}^{\beta}(x)&=(x)_a(x+\beta)_a
\bigl((x+a+\beta)_{b-a}-(x+a)_{b-a}\bigr).
\end{aligned}
\]

We obtain the exact finite identity

\[
Q_n^{\alpha,\beta}(x)=
\sum_{\substack{j\ge1,\ 0\le a<b\\a+b+j=n}}
\frac{n!(\alpha)_j}{j!a!b!}
M_{a,b}^{(j)}B_{a,b}^{\beta}(x).
\tag{4}
\]

Every \(M_{a,b}^{(j)}\) is nonnegative. Indeed, if \(f_af_{b+j}=0\),
this is immediate. Otherwise all intermediate terms are positive, and
the nonincreasing successive ratios of a log-concave sequence give

\[
\frac{f_{a+j}}{f_a}
=\prod_{t=1}^j\frac{f_{a+t}}{f_{a+t-1}}
\ \ge\ \prod_{t=1}^j\frac{f_{b+t}}{f_{b+t-1}}
=\frac{f_{b+j}}{f_b}.
\]

Equivalently, \(M_{a,b}^{(j)}\) is the Toeplitz minor with row indices
\(0,j\) and column indices \(a+j,b+j\).

The polynomial \(B_{a,b}^{\beta}\) has nonnegative coefficients. To see
this, put \(d=b-a\ge1\). Each factor in \((x+a+\beta)_d\) has a strictly
larger nonnegative constant term than the corresponding factor in
\((x+a)_d\). Their difference has strictly positive coefficients in
every degree \(0,\ldots,d-1\), and its leading degree-\(d\) coefficient
cancels. One may also use the telescoping product identity

\[
\begin{aligned}
&\prod_{t=0}^{d-1}(x+a+\beta+t)-\prod_{t=0}^{d-1}(x+a+t)\\
&\quad=\beta\sum_{u=0}^{d-1}
\prod_{t=0}^{u-1}(x+a+\beta+t)\prod_{t=u+1}^{d-1}(x+a+t).
\end{aligned}
\tag{5}
\]

Multiplication by \((x)_a(x+\beta)_a\) preserves coefficient
nonnegativity. Moreover,

\[
\deg B_{a,b}^{\beta}=a+b-1=n-j-1\le n-2.
\]

This proves the first assertion.

Under (3), all interior terms \(f_1,\ldots,f_{n-1}\) are positive. If
both endpoints are positive, strictly decreasing successive ratios imply
\(f_1/f_0>f_n/f_{n-1}\). If an endpoint is zero, the same strict
inequality after cross multiplication is immediate. Consequently

\[
M_{0,n-1}^{(1)}=f_1f_{n-1}-f_0f_n>0.
\]

The corresponding term of (4) is

\[
n\alpha\bigl(f_1f_{n-1}-f_0f_n\bigr)
\bigl((x+\beta)_{n-1}-(x)_{n-1}\bigr).
\tag{6}
\]

It has a strictly positive coefficient in every degree \(0,\ldots,n-2\).
Every other term has nonnegative coefficients. The conclusion follows.
\(\square\)

The proof does not infer coefficient positivity from stability. Section
4 shows that the proposed stability implication is false.

\Needspace{9\baselineskip}\section{3. Higher determinants: an expansion in Toeplitz
minors}\label{higher-determinants-an-expansion-in-toeplitz-minors}

Fix \(r\ge2\) and put \(R=\binom r2\). For a strictly increasing tuple
of nonnegative integers \(\ell=(\ell_0,\ldots,\ell_{r-1})\), write

\[
\begin{aligned}
V(\ell)&=\prod_{0\le i<j<r}(\ell_j-\ell_i),\\
M_f(\ell)&=\det[f_{\ell_j+r-1-i}]_{i,j=0}^{r-1},\\
B_\ell(x)&=\prod_{j=0}^{r-1}(x+j)_{\ell_j-j}.
\end{aligned}
\tag{7}
\]

The last formula is a polynomial because \(\ell_j\ge j\). The matrix
defining \(M_f\) is a Toeplitz submatrix with rows \(0,\ldots,r-1\) and
columns \(\ell_0+r-1,\ldots,\ell_{r-1}+r-1\).

\Needspace{5\baselineskip}\textbf{Theorem 2.} For every input sequence and every \(n\ge0\),

\[
P_n^r(x)=n!\sum_{\substack{0\le\ell_0<\cdots<\ell_{r-1}\\
\ell_0+\cdots+\ell_{r-1}=n-R}}
\frac{V(\ell)}{\prod_{j=0}^{r-1}\ell_j!}
M_f(\ell)B_\ell(x).
\tag{8}
\]

An empty sum is zero. In particular, \(P_n^r=0\) if \(n<r(r-1)\).

Suppose now that \(f\) is \(PF_r\), let \(n\ge r(r-1)\), and put
\(d=n-r(r-1)\).

\begin{enumerate}
\def\labelenumi{\arabic{enumi}.}
\tightlist
\item
  Every coefficient of \(P_n^r\) is nonnegative.
\item
  The output is nonzero if and only if some minor \(M_f(\ell)\) in (8)
  is positive. In that case its degree is exactly \(d\).
\item
  Every coefficient in degrees \(0,\ldots,d\) is strictly positive if
  and only if some contributing tuple with \(\ell_0=0\) has
  \(M_f(\ell)>0\).
\item
  If the output is nonzero and the condition in part 3 fails, then
  \(d\ge1\), its constant coefficient is zero, and all its coefficients
  in degrees \(1,\ldots,d\) are strictly positive. Its zero at the
  origin is simple.
\end{enumerate}

\textbf{Proof of the identity.} Let \(\Delta H(x)=H(x+1)-H(x)\).
Applying successive finite differences to the columns and then to the
rows of the matrix in (2) uses triangular operations with diagonal
entries one. The resulting entry in position \((i,j)\) is

\[
(-1)^i\Delta^{i+j}\Phi_f(x-i,z)
=(-1)^iz^{i+j}\sum_{k\ge0}
f_{k+i+j}(x+j)_k\frac{z^k}{k!}.
\]

Here we used \(\Delta^q(x)_k=k!\,(x+q)_{k-q}/(k-q)!\) for \(k\ge q\),
and zero otherwise. Taking the row and column factors out of the
determinant gives

\[
\det[\Phi_f(x+j-i,z)]
=(-1)^Rz^{2R}\det[H_{ij}(x,z)],
\tag{9}
\]

where

\[
H_{ij}(x,z)=\sum_{\ell\ge j}
f_{i+\ell}(x+j)_{\ell-j}\frac{z^{\ell-j}}{(\ell-j)!}.
\]

Factor this matrix as \(H=AB\), with

\[
A_{i\ell}=f_{i+\ell},\qquad
B_{\ell j}=\begin{cases}
(x+j)_{\ell-j}z^{\ell-j}/(\ell-j)!,&\ell\ge j,\\
0,&\ell<j.
\end{cases}
\]

Cauchy-\/-Binet applies coefficient by coefficient: truncating the
intermediate index makes the sum finite, and each fixed coefficient
stabilizes. For an increasing tuple \(\ell\), reversing the rows of its
\(A\) minor gives

\[
\det[f_{i+\ell_j}]_{i,j=0}^{r-1}=(-1)^R M_f(\ell).
\]

For generic \(x\) and nonzero \(z\), its \(B\) minor can be evaluated by
factoring its rows and columns:

\[
\det[B_{\ell_i,j}]_{i,j=0}^{r-1}
=z^{\sum_i\ell_i-R}
\frac{\prod_i(x)_{\ell_i}}{\prod_j(x)_j\prod_i\ell_i!}
\det[\ell_i^{\underline j}]_{i,j=0}^{r-1}.
\]

Since the falling-factorial polynomials \(u^{\underline j}\) are monic
of successive degrees, their determinant is \(V(\ell)\). Also

\[
\frac{\prod_i(x)_{\ell_i}}{\prod_j(x)_j}
=\prod_{j=0}^{r-1}(x+j)_{\ell_j-j}=B_\ell(x).
\]

These identities extend polynomially to all \(x\). Substitution into (9)
cancels the two signs \((-1)^R\) and gives powers
\(z^{\sum_i\ell_i+R}\). Extracting \(n![z^n]\) proves (8).

\textbf{Proof of the positivity statements.} Under \(PF_r\), all
\(M_f(\ell)\ge0\). Every other scalar in (8) is positive. Each
\(B_\ell\) is monic, has nonnegative coefficients, and has degree

\[
\sum_j(\ell_j-j)=n-2R=d.
\]

Thus a positive minor gives a strictly positive leading coefficient,
while vanishing of all the minors gives the zero polynomial. This proves
parts 1 and 2.

If \(\ell_0=0\), every linear factor in \(B_\ell\) has positive constant
term, so all coefficients through its degree are positive. If
\(\ell_0>0\), exactly one linear factor is \(x\), and every other linear
factor has positive constant term. Such a summand has zero constant
coefficient and strictly positive coefficients in degrees
\(1,\ldots,d\). The nonnegative weights in (8) now prove parts 3 and 4.
\(\square\)

\Needspace{5\baselineskip}\textbf{Corollary 3 (fixed integer zeros).} If \(d=n-r(r-1)>0\), then
every \(P_n^r\), with no positivity assumption, is divisible by

\[
(x+r-1)_{\lceil d/r\rceil}.
\tag{10}
\]

\textbf{Proof.} The integers \(q_j=\ell_j-j\) in each summand are
nonnegative and nondecreasing, and sum to \(d\). Hence
\(q_{r-1}\ge\lceil d/r\rceil\). The last factor in \(B_\ell\) supplies
(10). \(\square\)

\subsection{3.1. The necessary qualification to Conjecture
4}\label{the-necessary-qualification-to-conjecture-4}

The nonstrict \(PF_r\) hypothesis alone does not force the asserted
degree or strict coefficient positivity. For example, take \(r=2\) and
\(f_k=1\) for \(0\le k\le n\), with zero continuation. This sequence is
\(PF_2\) and has positive endpoints, but Vandermonde\textquotesingle s
identity gives

\[
P_n^2(x)=(2x)_n-(2x)_n=0.
\tag{11}
\]

Even nonzero outputs need not have positive constant coefficient. For
\(n=4\), the \(PF_\infty\) sequence \((1,2,1,0,0)\) gives

\[
P_4^2(x)=12x(x+1).
\tag{12}
\]

Thus the literal assertion in Conjecture 4 requires a nondegeneracy
qualification. Theorem 2 supplies its exact replacement. In particular,
when at least one specified minor with \(\ell_0=0\) is positive, the
degree and all the coefficient inequalities conjectured there follow in
every order \(r\).

\Needspace{9\baselineskip}\section{4. A single positive input disproves Conjectures 2 and
5}\label{a-single-positive-input-disproves-conjectures-2-and-5}

We give the input and a negative Hurwitz determinant explicitly.
Numerical roots are unnecessary for the disproof.

For a degree-eight polynomial

\[
T(x)=a_0x^8+a_1x^7+\cdots+a_8,\qquad a_0>0,
\]

define

\[
\Delta_k(T)=\det[a_{2j-i+1}]_{i,j=0}^{k-1},
\qquad a_t=0\ \text{for }t\notin\{0,\ldots,8\}.
\tag{13}
\]

The Hurwitz criterion {[}H95{]} implies that a Hurwitz stable polynomial
has \(\Delta_k>0\) for all \(k\). In fact, a negative Hurwitz
determinant forces a zero in the open right half-plane: if all zeros had
nonpositive real parts, then \(T(x+\varepsilon)\) would be Hurwitz
stable for every \(\varepsilon>0\), and continuity would force every
\(\Delta_k(T)\ge0\).

\Needspace{5\baselineskip}\textbf{Theorem 4.} There is a positive strictly log-concave finite
sequence for which \(Q_{18}^{1,1}\) and \(P_{18}^2\) both have zeros in
the open right half-plane.

\textbf{Proof.} Define

\[
h_k=k(18-k),\qquad 0\le k\le18,
\tag{14}
\]

and take

\[
f_0=f_{18}=1,\qquad
f_k=10001+h_k\quad(1\le k\le17).
\tag{15}
\]

All terms are positive and the sequence is symmetric. For
\(2\le k\le16\), put \(t=18-2k\). Since \(h_{k-1}=h_k-t-1\) and
\(h_{k+1}=h_k+t-1\),

\[
f_k^2-f_{k-1}f_{k+1}=2f_k-1+t^2>0.
\tag{16}
\]

At \(k=1\), the gap is \(10018^2-10033>0\), and the case \(k=17\)
follows by symmetry. This proves strict log-concavity; its zero
continuation is \(PF_2\).

Direct expansion of (2), or (8), gives

\[
P_{18}^2(x)=88128\,(x+1)_8\,R(x),
\tag{17}
\]

where

\[
\begin{aligned}
R(x)={}&4930033x^8+217997860x^7+4622212322x^6\\
&+66604657560x^5+741214112737x^4\\
&+6223968331700x^3+35647227456908x^2\\
&+120787851904080x+180829229981400.
\end{aligned}
\tag{18}
\]

Consequently \(Q_{18}^{1,1}(x)=88128\,(x+2)_8\,T(x)\), where
\(T(x)=R(x+1)\) is

\[
\begin{aligned}
T(x)={}&4930033x^8+257438124x^7+6286238266x^6\\
&+99191968400x^5+1151545612777x^4\\
&+9955221611636x^3+59506512884844x^2\\
&+214081390250960x+344300941584600.
\end{aligned}
\tag{19}
\]

Its seventh Hurwitz determinant is negative. A conveniently line-broken
exact certificate is

\[
\Delta_7(T)=-2^{16}3^7 5^5 7^3\,11\,13^3\,N<0,
\tag{20}
\]

where the positive integer \(N\) is specified without approximation by

\[
\begin{aligned}
N={}&378576120131718481193785864\cdot10^{33}\\
&+480712644384882588463864139\cdot10^6+262343.
\end{aligned}
\tag{21}
\]

Formula (13) and the nine coefficients in (19) suffice to check (20) by
integer arithmetic. The supplementary script does so from the original
input and also verifies (17) directly.

It follows that \(T\) has a zero \(\zeta\) with
\(\operatorname{Re}\zeta>0\). This is a zero of \(Q_{18}^{1,1}\),
disproving Conjecture 2 with its strict hypotheses. Since
\(R(\zeta+1)=T(\zeta)=0\), the polynomial \(P_{18}^2\) has a zero with
real part greater than one. The input is \(PF_2\), so this also
disproves Conjecture 5. \(\square\)

For orientation only, the offending conjugate pair of \(T\) is
approximately

\[
0.0363149803939\ \mathbin{\pm}\ 10.8821345313105\,i.
\tag{22}
\]

The exact determinant, rather than these approximations, is the
certificate. The full Hurwitz sign sequence is \((+,+,+,+,+,+,-,-)\);
the verification file records the exact integers.

\Needspace{9\baselineskip}\section{5. A binomial input disproves Conjecture
6}\label{a-binomial-input-disproves-conjecture-6}

\Needspace{5\baselineskip}\textbf{Theorem 5.} For \(r=3\), \(n=11\), and

\[
f_k=\binom{11}{k}\qquad(0\le k\le11),
\tag{23}
\]

the sequence is \(PF_\infty\) with positive endpoints, but \(P_{11}^3\)
has nonreal zeros.

\textbf{Proof.} The generating polynomial of the input is
\((1+z)^{11}\). One can verify \(PF_\infty\) directly: the Toeplitz
matrix of \((1,1)\) is totally nonnegative, and multiplying such
matrices corresponds to convolving their sequences. Cauchy-\/-Binet
preserves nonnegative minors. Eleven factors give (23). Thus no endpoint
or nonstrictness exception is involved.

Evaluation of (8) gives

\[
\begin{aligned}
P_{11}^3(x)&=15697281600\,(x+2)(x+3)C(x),\\
C(x)&=883x^3+5349x^2+9926x+5880.
\end{aligned}
\tag{24}
\]

This identity has a short finite certificate. Here \(R=3\), and the
increasing tuples of nonnegative integers with sum \(11-R=8\) are

\[
(0,1,7),\ (0,2,6),\ (0,3,5),\ (1,2,5),\ (1,3,4).
\tag{25}
\]

For additional transparency, their exact minors and weights
\(w_\ell=11!V(\ell)M_f(\ell)/\prod_j\ell_j!\) are

\[
\begin{array}{c|r|r}
\ell&M_f(\ell)&w_\ell\\\hline
(0,1,7)&330330&109880971200\\
(0,2,6)&1698840&2260408550400\\
(0,3,5)&2548260&4238266032000\\
(1,2,5)&2076360&4144082342400\\
(1,3,4)&1868724&3108061756800
\end{array}
\]

Thus \(P_{11}^3=\sum_\ell w_\ell B_\ell\), with the bases specified in
(7), gives (24) by five polynomial multiplications and additions.

For a real cubic \(ax^3+bx^2+cx+d\), the discriminant is

\[
b^2c^2-4ac^3-4b^3d-27a^2d^2+18abcd.
\]

For \(C\) it equals

\[
\operatorname{disc}(C)=-623476452956<0.
\tag{26}
\]

A real cubic with negative discriminant has one real zero and a nonreal
conjugate pair. Hence (24) disproves Conjecture 6 with all its input
conditions satisfied. \(\square\)

The nonreal pair is approximately
\(-1.43674838535\pm0.164358883618\,i\). Its negative real part explains
why this counterexample alone would not disprove Hurwitz stability;
Theorem 4 gives the separate stability obstruction.

\Needspace{9\baselineskip}\section{6. Relation to Conjecture 3 and
scope}\label{relation-to-conjecture-3-and-scope}

Theorem 5 has determinant order three. It does not contradict the
order-two real-rootedness result in the separate manuscript {[}Z26{]}.
That manuscript proves the conclusion of Karp\textquotesingle s
Conjecture 3 for every \(n\ge3\) when the finite \(PF_\infty\) input
satisfies \(f_0f_n>0\). If zero endpoints are admitted, its proved
conclusion is that the output is either zero or has only real
nonpositive zeros. Example (12) prevents a universal strictly negative
conclusion under that wider convention.

Theorem 4 uses a \(PF_2\) input; it makes no \(PF_\infty\) claim. Its
nonreal right-half-plane zeros therefore also do not contradict the
positive-endpoint result for Conjecture 3.

The results here settle the remaining numbered assertions as formulated:
Conjecture 1 holds; Conjectures 2, 5, and 6 fail; the unqualified
Conjecture 4 fails and has the exact replacement in Theorem 2. They do
not classify every stronger hypothesis under which the failed
zero-location assertions might hold.

\Needspace{9\baselineskip}\section{7. Exact verification and
provenance}\label{exact-verification-and-provenance}

The supplementary program \texttt{checks/verify.py} uses exact integers
and rational polynomials. It independently expands the defining
order-two expression for (15), verifies its factorization, and evaluates
the Hurwitz determinants. It checks the five-term expansion (25), the
cubic discriminant, and the order-three example against the original
determinant at thirteen distinct integer values of \(x\). Both
expressions in that last comparison have degree at most eleven, so those
evaluations provide an independent polynomial-identity certificate. It
also checks symbolic low-degree instances of the arbitrary-shift formula
and additional instances of (8).

No numerical root computation is used to certify a counterexample.
Optional approximations are labeled as such. Finite checks support the
algebraic proofs and do not replace Theorems 1 and 2 for arbitrary \(n\)
and \(r\).

This manuscript was prepared with OpenAI Codex under the
author\textquotesingle s instructions. The same assistant developed the
proofs and performed the computational and mathematical checks. No
independent human review, journal acceptance, or priority determination
is claimed. This manuscript has its own sources and verification files
and is not a revision of {[}Z26{]}.

\Needspace{9\baselineskip}\section{References}\label{references}

{[}H95{]} A. Hurwitz. Ueber die Bedingungen, unter welchen eine
Gleichung nur Wurzeln mit negativen reellen Theilen besitzt.
\emph{Mathematische Annalen} \textbf{46} (1895), 273-284.
\href{https://doi.org/10.1007/BF01446812}{DOI}.
\href{https://gdz.sub.uni-goettingen.de/download/pdf/PPN235181684_0046/LOG_0026.pdf}{Original
article scan}.

{[}K12{]} D. Karp. Positivity of Toeplitz determinants formed by rising
factorial series and properties of related polynomials. \emph{Zap.
Nauchn. Sem. POMI} \textbf{404} (2012), 184-198; English translation,
\emph{Journal of Mathematical Sciences} \textbf{193} (2013), 106-114.
\href{https://arxiv.org/abs/1203.1482}{arXiv:1203.1482}.
\href{https://doi.org/10.1007/s10958-013-1438-y}{DOI}.

{[}Z26{]} H. Zweiman. Negative real zeros of a rising-factorial
transform. Research preprint, version 1.0, September 22, 2026.
\href{https://github.com/Hwzw/algebra-number-theory-drafts/tree/main/papers/P56-karp-rising-factorial-zeros}{Separate
manuscript}. Prepared with OpenAI Codex; not peer reviewed.

\end{document}
